{\large\textbf{Physics of Weighing Scales}}

Various scales for measuring the mass of objects are used in daily life. This
question is about the physical principles related to the beam balance, the
Roberval balance. While these balances look similar, they have slightly
different structures and behave differently.

We assume that small friction at the pivots allows the balance to eventually
come to rest. However, this friction is sufficiently small that it does not
affect the equilibrium angle determined from torque balance. Therefore, friction
and air resistance may be neglected in the calculations.

\textbf{A. Sensitivity of the Beam Balance (2.5 pts)}

\begin{center}\includegraphics[width=0.45\linewidth]{figures/APhO_2026_Q1_fig1.jpg}\end{center}

\begin{center}Fig.1\end{center}

A beam balance consists of a beam(lever arm) that rotates about a fixed
axis(pivot or fulcrum) and two pans of equal mass suspended from each side of
the beam. If the masses placed on the pans differ, the beam tilts toward the
heavier side to reach equilibrium.

During the motion of the beam, the suspended pans may swing. Although the force
exerted by the system consisting of the pan and the object on the beam may vary
over time due to this swinging, we approximate the force as the total weight of
the pan and the object, neglecting the swinging effect.

If the beam tilts at a large angle for even a small mass difference, the scale
is considered sensitive. Part A of the question examines the issue of
sensitivity.

The beam is assumed to be a flat sheet with negligible thickness. Let $O$ be the
fixed point, and $L$ and $R$ be the points where the left and right pans are
suspended, respectively. The center of mass of the beam coincides with point
$O$ , as in fig.2. The axis of rotation passes through $O$ and is perpendicular
to the beam. The physical parameters and variables that may be related to the
beam balance and its sensitivity are as follows.

\begin{itemize}
\item $b$: the vertical distance between $O$ and the line connecting $L$ and $R$
\item $l$: the horizontal distance from the perpendicular bisector passing
through $O$ to points $L$ and $R$
\item $g$: gravitational acceleration
\item $M$: Mass of the beam
\item $m_1$: Total mass of the left pan and its load
\item $m_2$: Total mass of the right pan and its load
\end{itemize}

When $m_1 > m_2$, the beam tilts counter-clockwise by an angle $\theta_0$ to
reach equilibrium.

\begin{center}\includegraphics[width=0.71\linewidth]{figures/APhO_2026_Q1_fig2.png}\end{center}

\begin{center}Fig. 2.\end{center}

\textbf{A.1}\quad When the beam is tilted by an angle $\theta$
counter-clockwise from the horizontal, find the magnitude of the torque about O
exerted by the left pan and its load, taking the counter-clockwise direction as
positive. \hfill 0.3 pt

\textbf{A.2}\quad When the beam is tilted by an angle $\theta$
counter-clockwise from the horizontal, find the torque exerted by the right pan
and its load (total mass $m_2$) that tends to rotate the beam clockwise.
\hfill 0.3 pt

\textbf{A.3}\quad Express the tilt angle $\theta_0$ at equilibrium in terms of
the given variables and parameters. \hfill 0.4 pt

\textbf{A.4}\quad To make the scale more sensitive (a larger $\theta_0$ for a
small mass difference), which of the following conditions for $b$ and $l$ is
correct? (Selecting an incorrect option will result in a 0.1-point deduction.)
\hfill 0.3 pt\\
\hspace*{1.5em}1. Larger $l$ or larger $b$ leads to a larger $\mid \theta_0 \mid$.\\
\hspace*{1.5em}2. Smaller $l$ or smaller $b$ leads to a larger $\mid \theta_0 \mid$.\\
\hspace*{1.5em}3. Larger $l$ or smaller $b$ leads to a larger $\mid \theta_0 \mid$.\\
\hspace*{1.5em}4. Smaller $l$ or larger $b$ leads to a larger $\mid \theta_0 \mid$.

The beam of a commercially available beam balance is often made so that the
rotation axis (pivot point $O$) is higher than the center of mass (CM) of the
beam. However, making the beam this way reduces the sensitivity of the beam
balance. To solve this problem and design a more sensitive scale, we intend to
change the structure of the beam. As a candidate, the beam is designed by
modifying it so that the pivot point ($O$) of the beam is below the center of
mass (CM) of the beam as shown in the fig. 3. Let the pivot point of the beam
positioned at a distance $d$ underneath the center of mass. The beam is assumed
to be a flat sheet with negligible thickness. The meanings of
$M, L, R, b, l, m_1, m_2, g$ for the scale are the same as in the previous
problem.

\begin{center}\includegraphics[width=0.71\linewidth]{figures/APhO_2026_Q1_fig3.png}\end{center}

\begin{center}Fig.3\end{center}

\textbf{A.5}\quad When the beam tilts by an angle $\theta_1 (< \pi/2)$ from the
horizontal to reach equilibrium, express the tilt angle $\theta_1$ in terms of
the given variables and parameters. \hfill 0.8 pt

\textbf{A.6}\quad Obtain the condition under which the beam reaches a stable
equilibrium angle $\theta_1 (< \pi/2)$. Express the condition as an inequality
independent of $\theta_1$. \hfill 0.4 pt

\textbf{B. Basic Model of Roberval Balance (3.6 pts)}

\begin{center}\includegraphics[width=0.45\linewidth]{figures/APhO_2026_Q1_fig4.jpg}\end{center}

\begin{center}Fig.4\end{center}

The Roberval balance uses a parallel-linkage structure, where the pans are
connected to two horizontal beams (upper and lower). These two beams are joined
to the pans by pivots, which act like hinges. This special connection allows
each pan of two pivots to stay perfectly vertical even when the beams tilt
(Fig.4). As the beams rotate, the pans move together in a synchronized way. A
unique feature of this design is that the balance depends only on the total
mass on each side; it does not matter where you place the weights on the pans.
The physical parameters, variables and notations that may be related to the
beam balance is as follows (Fig.5).

\begin{itemize}
\item $O, O'$: Fixed pivots for the two horizontal beams
\item $I_1$: The moment of inertia of the upper beam about its axis of rotation
\item $I_2$: The moment of inertia of the lower beam about its axis of rotation
\item $l$: The distance from the central pivot to the pan suspension point
\item $x_L, x_R$: Horizontal offsets of the weights from the center of the left
and right pans, respectively.
\item $m$: Mass of each pan
\item $m_L, m_R$:Mass of the load placed on the left and right pans,
respectively. ($m_L \geq m_R$)
\item $g$: Gravitational acceleration.
\end{itemize}

Assume that the Center of mass (CM) of each beam coincides with its pivot and
that pivots of pans and pivot of the beam lays on one line.

\begin{center}\includegraphics[width=0.85\linewidth]{figures/APhO_2026_Q1_fig5.png}\end{center}

\begin{center}Fig.5\end{center}

\textbf{B.1}\quad Calculate the total potential energy of the system
$U(\theta)$, when the beam is tilted counter-clockwise by an angle $\theta$ from
the horizontal ($m_L \geq m_R$). Define the potential energy $U$ to be zero at
the initial horizontal position. \hfill 0.3 pt

\textbf{B.2}\quad Express the total kinetic energy of the system in terms of the
given variables and parameters and the angular velocity $\dot{\theta}$.
\hfill 0.5 pt

\textbf{B.3}\quad Obtain the second-order differential equation governing the
rotation angle $\theta$. \hfill 0.6 pt

The angular acceleration $\ddot{\theta}$ at the instant the beam is released
from the horizontal position is as follows:
\[
\ddot{\theta} = \frac{(m_L - m_R)gl}{I_1 + I_2 + (2m + m_L + m_R)l^2}
\]

\textbf{B.4}\quad At the instant with zero initial velocity, let
$T_{L1}, T_{L2}$ be the magnitudes of the vertical components of forces acting
between the left pan and the upper and lower beams, respectively. and
Similarly, let $T_{R1}, T_{R2}$ be the magnitudes of the vertical components of
the forces for the right pan. Calculate the values of $(T_{L2} + T_{R1})$ in
terms of the given variables and parameters. \hfill 1 pt

\textbf{B.5}\quad Assuming that all components of the balance, including the
beams and pans, are rigid bodies, determine whether each of the following
forces can be calculated at the moment of release. (Answer with Yes, No, or
blank for each. A penalty of 0.1 points will be applied for each incorrect
answer.) \hfill 0.6 pt\\
\hspace*{1.5em}1. $T_{R1}$\\
\hspace*{1.5em}2. Vertical component of force exerted by the central pivot on
the upper beam\\
Set up all the necessary relational equations required to solve this problem.
Note that you are only required to provide the forms of the equations; explicit
final calculations are not necessary.

\textbf{B.6}\quad Let $M_T$ be the mass of the balance without any weights.
Weights of mass $m_L$ and $m_R$ ($m_L > m_R$) are placed on the left and right
pans, respectively. The beam is initially held horizontal by hand and then
released. Find the normal force $N$ exerted by the floor on the balance
immediately after release. \hfill 0.6 pt

\textbf{C. Practical Model of Roberval Balance (3.9 pts)}

In the basic model of Roberval balance discussed in Part B, a mass imbalance
causes continuous angular acceleration, making it impossible to determine a
static equilibrium angle. In contrast, a practical Roberval balance reaches a
stable equilibrium at a specific tilt angle depending on the mass difference.
In Part C, we analyze the physical structure of such practical Roberval
balances.

To calculate the equilibrium angle as a function of the mass difference,
consider the following variables and parameters:

\begin{itemize}
\item Upper Beam: The pivot point (fixed axis of the beam) is located at a
vertical distance $d$ directly above the beam's center of mass. The beam has a
mass $M$ and a moment of inertia $I_1$ about its pivot(fixed axis).
\item Lower Beam: The pivot point(fixed axis of the beam) coincides with the
beam's center of mass. The beam has a moment of inertia $I_2$ about its
pivot(fixed axis).
\item $m_1, m_2$ ($m_1 \geq m_2$) : The combined mass of the pan and any objects
placed upon it for the left and right sides, respectively. (Please note the
difference in notation from part B.)
\item $l$: The horizontal distance from the perpendicular bisector passing
through the central pivot to the pan suspension point
\item $g$: Gravitational acceleration.
\end{itemize}

It is assumed that the weights remain stationary relative to the pans and move
in unison with them and that pivots of pans and pivot of the beam lays on one
line.

\textbf{C.1}\quad With two weights of different masses placed on the
pans($m_1 > m_2$), the beam is initially held in a horizontal position and then
released from rest. In the situation, the second-order differential equation of
motion for the tilt angle $\theta$ takes the form:
$A\ddot{\theta} = B\cos\theta + C\sin\theta$ . Determine the coefficients
$A, B$ and $C$ in terms of the given variables and parameters. Set
$\theta = 0$ at the horizontal position. \hfill 1.4 pt

\textbf{C.2}\quad When the balance is at its equilibrium state
($\theta = \theta_0$), a slight disturbance causes the beams and pans to
oscillate about the equilibrium angle. To analyze this small oscillation, we
define a new variable $\eta = \theta - \theta_0$. By approximating the equation
of motion obtained in Part C.1, derive the governing equation for $\eta$ in
terms of the given variables and parameters. The answer must not include
$\theta_0$. \hfill 1.9 pt

\textbf{C.3}\quad If the total mass $m_1 + m_2$ is constant, determine how the
mass should be distributed between the pans to maximize the period of small
oscillations. Calculate the period of small oscillations in the limit where
$m_1 = m_2 = 0$. \hfill 0.6 pt
